\documentclass[11pt,a4paper]{article}
\usepackage[margin=1in]{geometry}
\usepackage{booktabs}
\usepackage{longtable}
\usepackage{graphicx}
\usepackage{hyperref}
\usepackage{listings}
\usepackage{xcolor}
\usepackage{enumitem}
\usepackage{caption}
\usepackage{amsmath}
\usepackage{amssymb}
\usepackage{tcolorbox}
\tcbuselibrary{breakable}

\hypersetup{
    colorlinks=true,
    linkcolor=blue,
    filecolor=magenta,
    urlcolor=cyan,
}

\tcbset{
    promptbox/.style={
        colback=blue!5!white,
        colframe=blue!75!black,
        title={#1},
        fonttitle=\bfseries,
        before skip=6pt,
        after skip=6pt,
        breakable,
    }
}

% Long unbreakable \texttt{} paths are the main source of overfull boxes in
% this report; a generous emergencystretch plus sloppy paragraph breaking lets
% TeX loosen those lines instead of running into the margin.
\setlength{\emergencystretch}{3em}
\sloppy

\title{Performance Report 24: Vulkan Engine \\
\large The Fullscreen Tail --- Composite, Sky, and Brush Passes Every Pixel Pays After the Scene Renders}
\author{}
\date{\today}

\begin{document}

\maketitle

\begin{abstract}
Reports~17--23 audited everything up to the scene targets: ray tracing, the
water machinery and its noise, the solid/water raster core, vegetation,
submit counts, buffer pools, texture arrays, the GPU mixer, and the load
path. This report audits what happens \emph{after}: the fullscreen
post-process composite that assembles every frame, the sky that is shaded
twice (once fullscreen inside the solid pass, once as a cached equirect
fetched by the composite), and the brush early/backface/overlay tail that
runs three full-resolution passes per frame whether or not any brush is
visible.

The composite executes an unconditional prologue per pixel --- scene color
and depth fetches, an inverse-view-projection reconstruct with a normalize,
\texttt{atan}+\texttt{acos} equirect mapping, a sky fetch, plus vegetation,
brush, and two debug depth fetches --- before testing a single alpha
(\texttt{shaders/postprocess.frag:104--114,142,233,250,258}). The sky
fullscreen draw sits \emph{before} the \texttt{renderSolid} gate
(\texttt{MyApp.cpp:1989--1994} vs \texttt{:2000}), so it runs even with solid
rendering disabled, evaluating three \texttt{pow}s, a PCG star hash that is
multiplied by zero all day, and a \texttt{pow}-1200 sun flare per pixel
(\texttt{shaders/sky.frag:30,41--42,62--66,74}). The brush tail records three
clear-plus-draw full-resolution passes and a dedicated queue submit every
frame with no visibility or mesh-count gate (\texttt{MyApp.cpp:1668},
\texttt{BrushRenderer.cpp:214--345}), reusing tessellated solid shaders for
depth-only work. The equirect cache (report~22 H8) is sound in mechanism but
blind in signature: night colors, star intensity, and sky mode are not
compared (\texttt{MyApp.cpp:1761--1774} vs
\texttt{widgets/SkySettings.hpp}), so those edits bake a stale sky. Tail
targets commit \textbf{$\sim$95\,MiB} device-local at 1080p ($\sim$380\,MiB
at 4K) with no scaling.

It finds \textbf{15 issues} (3 critical, 5 high, 4 medium, 3 low). Each
carries the standard technique and a copy-pasteable prompt preserving
validation-clean execution, Synchronization2, and runtime feature detection.
Scope score \textbf{5.1/10}; folding the tail slice into the combined
estimate at 10\% weight gives \textbf{6.1/10}.
\end{abstract}

\tableofcontents
\newpage

% =====================================================================
\section{Executive Summary}
% =====================================================================

\subsection{Scope and Method}

Three areas, none previously reported on as subsystems: (1) the post-process
composite (\texttt{vulkan/renderer/PostProcessRenderer.cpp} + 268-line
\texttt{shaders/postprocess.frag}, 17 bindings, one fullscreen triangle per
frame on every preset); (2) the sky --- fullscreen gradient/grid pass inside
the solid color pass (\texttt{vulkan/renderer/SkyRenderer.cpp:81--119},
\texttt{shaders/sky.frag}), the 2048$\times$1024 cached equirect
(\texttt{:155--330}), and the per-frame \texttt{SkySphere} upload; (3) the
brush tail --- early front-depth/color passes plus the GREATER backface pass
(\texttt{vulkan/renderer/BrushRenderer.cpp:214--345},
\texttt{BrushBackFaceRenderer.cpp:120--197}) composited inside the
post-process. Ray tracing, water, raster core, vegetation, submits, pools,
arrays, mixer, and load path are referenced only at interfaces. Method
matches reports~22--23: full reads with file:line evidence,
\texttt{spirv-dis} on shipped modules, byte tallies from creation code. The
application was not executed; dynamic magnitudes are analytic estimates
from verified structure at $1920\times1080$ (2.07\,M pixels), 3 frames in
flight.

\subsection{Measured Budgets}

\begin{table}[h]
\centering
\small
\begin{tabular}{@{}lrrl@{}}
\toprule
\textbf{Module} & \textbf{Words} & \textbf{Instrs} & \textbf{Role} \\
\midrule
\texttt{main\_brush.frag.spv} & 22\,511 & 4\,809 & brush color (full \texttt{main.frag} $-$DBRUSH\_PASS) \\
\texttt{postprocess.frag.spv} & 4\,600 & 949 & fullscreen composite, 17 bindings \\
\texttt{sky.frag.spv} & 4\,159 & 846 & fullscreen sky gradient + stars + flare \\
\texttt{sky\_equirect.frag.spv} & 4\,082 & 822 & 2\,M-px cached bake, same sky math \\
\texttt{sky\_grid.frag.spv} & 3\,535 & 710 & grid-mode twin of the sky FS \\
\texttt{sky\_fullscreen.vert.spv} & 2\,515 & 515 & 3-vertex triangle + invVP reconstruct \\
\texttt{fullscreen.vert.spv} & 306 & 64 & plain 3-vertex triangle (composite) \\
\texttt{debug\_sdf.frag.spv} & 850 & 170 & SDF debug leg (fetched via composite) \\
\texttt{depth\_linearize.frag.spv} & 589 & 130 & debug-widget preview only \\
\texttt{depth\_only.frag.spv} & 218 & 50 & brush backface FS \\
\bottomrule
\end{tabular}
\caption{Disassembly word counts (\texttt{spirv-dis | wc -w}, the
reports~20--23 convention) and instruction counts (op lines). The brush
variant of the solid fragment shader is $5\times$ the composite and
$5.4\times$ the sky shader: the editing overlay shades with the full solid
recipe.}
\end{table}

Per-frame fixed structure: one composite fullscreen (2.07\,M px, every
preset), one sky fullscreen inside the solid pass, three brush
clear-plus-draw full-resolution passes plus a dedicated queue submit, one
UBO upload plus 17-entry signature check for the composite, one
\texttt{SkySphere} memcpy; the equirect costs 0 on a cache hit or 2.1\,M px
on a miss. Committed tail-side memory (\S\ref{sec:vram}):
\textbf{$\sim$95\,MiB device} at 1080p, $\sim$4$\times$ at 4K.

\subsection{Recommended Order}

Gate the brush tail on visibility/mesh count (C1) $\rightarrow$ gate the sky
fullscreen on \texttt{renderSolid} and day-factor the star hash (C2)
$\rightarrow$ defer the composite prologue past the opaque-solid test (C3)
$\rightarrow$ fix the equirect signature (H4) and de-tessellate brush depth
(H5) $\rightarrow$ shrink brush shading and clamp/skip the blur (H6, H7)
$\rightarrow$ tier the equirect (H8) $\rightarrow$ dirty-gate the uploads and
fetches (M9--M11) $\rightarrow$ scale the tail targets (M12) $\rightarrow$
hygiene (L13--L15).

% =====================================================================
\section{Cost Model}
% =====================================================================

\begin{equation}
\begin{aligned}
C = &\underbrace{\sum_{\mathrm{passes}} G_{\mathrm{pass}} V C_{\mathrm{vert}} + S C_{\mathrm{submit}} + L C_{\mathrm{bringup}} + G C_{\mathrm{sweep}}}_{\text{reports 20--23}} \\
&+ \underbrace{M_{\mathrm{tex}} C_{\mathrm{pressure}} + P\,C_{\mathrm{tail}} + B\,C_{\mathrm{brush}} + K\,C_{\mathrm{equirect}}}_{\text{this report}}
\end{aligned}
\end{equation}

with $P \approx 2.07$\,M pixels/frame at 1080p ($8.3$\,M at 4K),
$B = 1$ unconditional brush tail (three full-res passes + one queue submit),
and $K \in \{0, 1\}$ equirect bakes per frame (2.1\,M px on a miss).
Reports~20--23 priced geometry, submits, bring-up, generation sweeps, and
residency; this report prices the per-pixel tail, which scales with output
resolution rather than scene size --- the only term that quadruples at 4K
while everything else stays flat.

% =====================================================================
\section{Tail Residency}
\label{sec:vram}
% =====================================================================

Tallied from creation code at $1920\times1080$, 3 frames in flight
(\texttt{VulkanApp::MAX\_FRAMES\_IN\_FLIGHT = 3},
\texttt{BrushRenderer.hpp:42}, \texttt{SkyRenderer.hpp:73}).

\begin{table}[h]
\centering
\small
\begin{tabular}{@{}p{5.9cm}p{4.3cm}r@{}}
\toprule
\textbf{Resource} & \textbf{Sizing} & \textbf{Bytes} \\
\midrule
Brush color targets & $3 \times 1920\times1080$ RGBA8 & 23.7\,MiB dev \\
Brush front-depth targets & $3 \times 1920\times1080$ D32 & 23.7\,MiB dev \\
Brush backface-depth targets & $3 \times 1920\times1080$ D32 & 23.7\,MiB dev \\
Sky equirect targets & $3 \times 2048\times1024$ RGBA8 & 24.0\,MiB dev \\
Post-process UBO + descriptor buffers & WaterUBO + 3 slots & $<$0.1\,MiB host \\
\midrule
Total & & $\sim$95\,MiB dev \\
\bottomrule
\end{tabular}
\caption{Committed tail memory. None of it scales with scene content; all
of it scales with output resolution ($\sim$380\,MiB at 4K). Report~22's
pool tally did not include these targets.}
\end{table}

The brush targets are full-screen-resolution even though brush meshes are a
handful of chunks; the equirect is fixed $2048\times1024$ even on the
Minimal preset. Both are allocated per frame slot, unconditionally.

% =====================================================================
\section{Critical Issues}
% =====================================================================

\subsection{C1: The Brush Tail Runs Three Full-Res Passes Plus a Queue Submit Every Frame, Ungated}
\label{sec:c1}

\textbf{Severity: Critical.}

\texttt{MyApp.cpp:1668} records \texttt{BrushRenderer::recordEarlyPass} on
the dedicated brush-solid queue every frame with no visibility, mesh-count,
or settings gate around the block. That function
(\texttt{BrushRenderer.cpp:214--345}) performs three full-resolution
\texttt{BeginRendering} passes per frame: front depth (clear 1.0 + tessellated
depth draw, \texttt{:266--289}), backface depth via
\texttt{renderBackFacePass} (clear 0.0 GREATER + draw,
\texttt{BrushBackFaceRenderer.cpp:120--197}), and color (clear + draw,
\texttt{BrushRenderer.cpp:296--327}) --- each with viewport/scissor setup,
transitions in, and transitions to \texttt{SHADER\_READ\_ONLY} after
(\texttt{:329--344}). The brush cull inputs are likewise unconditional
(\texttt{MyApp.cpp} cull task always runs \texttt{acquireBuffers} +
\texttt{prepareCull} on the brush solid IR). With no brush mesh present the
draws are empty but the three clears, six transitions, the command buffer,
and the queue submit (with its timeline signal the composite waits on) all
still execute: $3 \times 2.07$\,M cleared pixels plus submit latency per
frame for an empty overlay.

\textit{Technique: gate on content, keep the graph.} Skip recording (and the
submit/signal) when the brush IRs hold zero meshes and no brush is armed;
the composite already treats missing brush input as transparent, so the
wait set can take the cull timeline directly on skipped frames, mirroring
the report-22 SDF-task elision (\texttt{MyApp.cpp:2275--2300}, which elides
the whole SDF task when disabled).

\begin{tcolorbox}[promptbox={Prompt --- C1: content-gated brush tail}]
Goal: an empty brush costs nothing per frame. Steps: (1) expose brush mesh
count / armed state from \texttt{BrushRenderer}; (2) skip the early-pass
record and the brush-solid submit when both IRs are empty and no brush
stroke is active, routing the composite wait to the already-signalled cull
timeline; (3) keep the targets' layouts stable across skipped frames so the
first stroke back needs no extra transition. Files:
\texttt{MyApp.cpp}, \texttt{vulkan/renderer/BrushRenderer.cpp,.hpp}.
Validation: identical pixels with and without brush, submit count per frame
down by one when idle, validation-clean.
\end{tcolorbox}

\subsection{C2: Sky Fullscreen Draws Before the renderSolid Gate and Shades Stars at Noon}
\label{sec:c2}

\textbf{Severity: Critical.}

The sky fullscreen draw (\texttt{MyApp.cpp:1989--1994}) sits before the
\texttt{if (settings.renderSolid)} gate (\texttt{:2000}), so disabling solid
rendering still pays the fullscreen sky: 2.07\,M pixels $\times$
\texttt{sky.frag} (4\,159 words / 846 instrs). That shader's per-pixel bill
is heavy and partly unconditional: three \texttt{pow} calls (sunset bias,
exponent bias, gradient falloff, \texttt{sky.frag:30,41--42}); a two-hash
PCG star seed (\texttt{:62--64}) whose result is multiplied by
\texttt{starMask} --- zero for the entire day (\texttt{:59,66}), so every
daytime pixel pays two hashes for speckles scaled to nothing; and a sun
flare \texttt{pow(max(sunDot,0), 800$\cdot$(1$-$sunElev$\cdot$0.5)}
(\texttt{:74}) with exponents up to $\sim$1200 evaluated per pixel even
looking away from the sun. Meanwhile the same sky is shaded a second time:
the composite reconstructs the view direction and fetches the equirect for
every background pixel (\texttt{postprocess.frag:107--114}) --- the
fullscreen pass and the equirect fetch overlap on exactly the pixels the
fullscreen pass just shaded.

\textit{Technique: gate the pass, factor the day.} Move the sky draw under
the \texttt{renderSolid} condition (or its own sky-visibility flag for
free-camera modes); branch the star block on \texttt{starMask > 0} so day
pixels skip both hashes; clamp the flare exponent and skip it for
\texttt{sunDot $\le$ 0} before the \texttt{pow}.

\begin{tcolorbox}[promptbox={Prompt --- C2: conditional sky fullscreen, day-factored stars}]
Goal: no sky ALU on pixels that cannot show it. Steps: (1) draw the sky
fullscreen only when solid rendering (or an explicit sky flag) is on;
(2) guard the \texttt{floatBitsToUint}+PCG pair with \texttt{starMask > 0};
(3) early-out the flare for \texttt{sunDot <= 0} and clamp the exponent.
Files: \texttt{MyApp.cpp}, \texttt{shaders/sky.frag}. Validation:
pixel-identical sky output, fullscreen sky cost zero with solid off,
validation-clean.
\end{tcolorbox}

\subsection{C3: The Composite Pays Its Full Prologue on Opaque-Solid Pixels}
\label{sec:c3}

\textbf{Severity: Critical.}

\texttt{postprocess.frag:main} runs the same prologue for every pixel
before any alpha is tested: scene color + depth fetches (\texttt{:104--105}),
clip$\rightarrow$world reconstruct through the full inverse view-projection
matvec plus divide plus \texttt{normalize} (\texttt{:107--111}),
\texttt{atan}+\texttt{acos}+divides equirect mapping plus a sky fetch
(\texttt{:113--114}) --- then line~119 reads \texttt{sceneColor.a}, which is
1.0 on the majority of pixels in any outdoors scene, meaning the sky color,
world position, and view direction just computed are discarded by the
\texttt{mix} on the very next line. After that, four more fetches are
unconditional: vegetation depth (\texttt{:142} else-branch, plus the 4-tap
\texttt{min} variant at \texttt{:134--143} when scaled), brush color
(\texttt{:233}), and both debug depths (\texttt{:250,258}) --- the
\texttt{$< 1.0$} tests happen \emph{after} the fetch, so cleared-to-1.0
targets (SDF/bbox passes elide when disabled, \texttt{MyApp.cpp:2275--2300},
leaving cleared targets) still cost fullscreen fetches every frame.

\textit{Technique: test first, reconstruct lazily.} Fetch scene color/depth,
branch on \texttt{isSolid == 1} to skip the world reconstruct + equirect +
sky fetch (compute them only for background and water/vegetation edge
pixels that need them); gate the veg/brush/debug fetches on CPU-provided
valid flags already known on the host (veg enabled, brush present this
frame, SDF/bbox enabled). Divergence cost is bounded: opaque-solid pixels
are the coherent majority.

\begin{tcolorbox}[promptbox={Prompt --- C3: lazy composite prologue}]
Goal: opaque pixels cost two fetches and a mix. Steps: (1) reorder
\texttt{main} to test \texttt{sceneColor.a} first and skip the
invVP/normalize/equirect/sky work when 1.0; (2) add UBO valid flags for
veg/brush/sdf/bbox (host-known) and skip those fetches when clear;
(3) keep the water/blur path untouched. Files:
\texttt{shaders/postprocess.frag},
\texttt{vulkan/renderer/PostProcessRenderer.cpp},
\texttt{vulkan/ubo/WaterUBO.hpp}. Validation: pixel-identical composite,
per-pixel fetch count on solid pixels down, validation-clean.
\end{tcolorbox}

% =====================================================================
\section{High-Severity Issues}
% =====================================================================

\subsection{H4: The Equirect Cache Signature Ignores Night, Stars, and Mode}
\label{sec:h4}

\textbf{Severity: High.}

Report~22 H8's cache compares a 12-float signature --- day horizon/zenith
(6), warmth, exponent, sunFlare (3), light direction (3)
(\texttt{MyApp.cpp:1761--1774}, exact \texttt{!=} compare at \texttt{:1769},
warmup 3 frames at \texttt{:276,1771--1772}). \texttt{SkySettings}
(\texttt{widgets/SkySettings.hpp}) carries five more inputs the equirect
shader reads: \texttt{nightHorizon}, \texttt{nightZenith},
\texttt{nightIntensity}, \texttt{starIntensity}, and \texttt{mode} (Gradient
vs Grid selects a different fullscreen pipeline on screen, while the cached
equirect keeps baking the old sky). Editing any of them at dusk leaves the
composite, water reflections, and solid360 sampling a stale 2\,M-px bake
until some unrelated day parameter changes. The compare is also exact float
equality, so a UI drag that lands back on the same value still invalidates.

\textit{Technique: sign everything the shader reads.} Extend the signature
to all \texttt{SkySettings} fields plus light elevation, or hash the
\texttt{SkyUniform} bytes the bake consumes; compare with a small epsilon.

\begin{tcolorbox}[promptbox={Prompt --- H4: complete equirect cache signature}]
Goal: every sky edit invalidates exactly once. Steps: (1) add night colors,
night/star intensities, and mode to the compared signature (or compare the
uploaded \texttt{SkyUniform} bytes); (2) use epsilon comparison for
drag-produced floats; (3) keep the 3-frame warmup. Files: \texttt{MyApp.cpp},
\texttt{widgets/SkySettings.hpp}. Validation: night/star/mode edits rebake
once, no-change frames still hit, validation-clean.
\end{tcolorbox}

\subsection{H5: Brush Depth Reuses the Tessellated Solid Pipeline for Depth-Only Work}
\label{sec:h5}

\textbf{Severity: High.}

Both brush depth passes run hull+domain shaders for depth-only output. The
front-depth pass draws through the solid deferred depth pipeline
(\texttt{SolidRenderer::drawDepthExternal},
\texttt{SolidRenderer.cpp:555--565}, binding
\texttt{activeDeferredDepthPipeline()}), and the backface pass builds its
own pipeline from \texttt{main.vert} + \texttt{main.tesc} +
\texttt{main.tese} + \texttt{depth\_only.frag} explicitly
(\texttt{BrushBackFaceRenderer.cpp:26--44}) --- tessellation factors
evaluated, patches subdivided, displaced vertices computed, all to write a
depth buffer at GREATER compare. The engine already owns no-tess vertex
twins for the main pass (report~21); neither brush depth path uses one.
Brush geometry is a few chunks, so this is vertex-rate waste rather than a
frame-rate cliff --- but it is the same HS/DS-per-depth-pixel pattern
report~21 removed from the main depth path, duplicated twice in the tail.

\textit{Technique: wire the no-tess twin.} Build the brush front/back depth
pipelines from the non-tessellated vertex shader (same bindings, no TCS/TES
stages); keep the tessellated path only where brush meshes opt into
displacement, if anywhere.

\begin{tcolorbox}[promptbox={Prompt --- H5: de-tessellated brush depth}]
Goal: brush depth costs vertices, not patches. Steps: (1) add no-TCS/TES
variants of the brush front-depth and backface pipelines; (2) select them
for brush IR draws; (3) confirm depth output identical. Files:
\texttt{vulkan/renderer/SolidRenderer.cpp},
\texttt{vulkan/renderer/BrushBackFaceRenderer.cpp}. Validation:
pixel-identical brush depth targets, HS/DS invocations zero on brush depth,
validation-clean.
\end{tcolorbox}

\subsection{H6: Brush Color Shades With the Full 22k-Word Solid Recipe}
\label{sec:h6}

\textbf{Severity: High.}

\texttt{shaders/main\_brush.frag.spv} is \texttt{shaders/main.frag} (182
lines plus heavy includes) compiled with \texttt{-DBRUSH\_PASS}
(\texttt{Makefile:146--150}) and still measures 22\,511 words / 4\,809
instructions --- $5\times$ the composite, $5.4\times$ the sky shader. Two
pipelines share it (blended deferred brush color and opaque brush overlay,
\texttt{SolidRenderer.cpp:444--497}), differing only in blend enable. The
\texttt{BRUSH\_PASS} guard trims PAINT mode, shadows, and set-1 declarations
(\texttt{shaders/main.frag:50,77}), but what remains is still the full solid
lighting/material machinery evaluated per brush pixel for an editing
gizmo that the composite then alpha-mixes over the scene
(\texttt{postprocess.frag:234--244}). An unlit or flat-shaded brush
fragment would be visually equivalent at these opacities
(\texttt{brushAlpha} default 0.5, \texttt{MyApp.cpp:3179--3185}).

\textit{Technique: a dedicated cheap brush FS.} Write a small brush fragment
shader (base color + simple lambert or unlit + constant alpha) behind the
same interface; keep the full-recipe path behind a quality flag if
pixel-parity is ever required.

\begin{tcolorbox}[promptbox={Prompt --- H6: cheap brush fragment shader}]
Goal: brush pixels cost a lambert term, not the solid recipe. Steps:
(1) add \texttt{shaders/brush\_flat.frag} (unlit/lambert, same outputs);
(2) point both brush pipelines at it; (3) A/B the overlay look at default
opacity. Files: \texttt{shaders/brush\_flat.frag} (new),
\texttt{vulkan/renderer/SolidRenderer.cpp}, \texttt{Makefile}.
Validation: overlay visually equivalent, brush-pass ALU down by an order of
magnitude, validation-clean.
\end{tcolorbox}

\subsection{H7: The 9-Tap Bilateral Water Blur Runs Unclamped exp() per Tap in the Composite}
\label{sec:h7}

\textbf{Severity: High.}

Report~19 H4 gated the body/column block on \texttt{waterBlurEnabled}
(\texttt{postprocess.frag:168,215}, \texttt{PostProcessRenderer.cpp:378--380},
\texttt{MyApp.cpp:3221--3225}) --- the residual inside the gate is still
expensive: per water pixel with \texttt{blurPx > 0.5}, one center fetch pair
plus 8 taps each costing a body fetch, a column fetch, and
\texttt{exp(-abs(dS - dC) * 0.5)} (\texttt{postprocess.frag:180--189}) ---
16 \texttt{textureLod} + 8 \texttt{exp} per pixel, with tap offsets scaled by
the unbounded material radius (\texttt{columnCenter.g}, \texttt{:171--173}),
so large radii sample sparsely across the body target and thrash the texture
cache exactly on the wide-water pixels the blur exists for. No mip is used
(all \texttt{Lod} 0.0); a 2-level mip pyramid on the body target would
approximate large radii with one fetch.

\textit{Technique: clamp the radius, mip the body.} Clamp
\texttt{blurPx} to a cache-sane maximum (spillover handled by the weight
normalization already present); generate 1--2 mips on the water body target
and select the level from the radius so wide blurs cost one filtered fetch
instead of 8 sparse ones.

\begin{tcolorbox}[promptbox={Prompt --- H7: bounded water-blur kernel}]
Goal: wide water costs fetches, not cache misses. Steps: (1) clamp
\texttt{blurPx} (e.g.\ $\le 8$\,px, retune per material); (2) add body-target
mips and level-select by radius; (3) keep the bilateral weight and the
reflection-sharpness invariant. Files: \texttt{shaders/postprocess.frag},
water body target creation. Validation: blurred-water A/B, water-pixel cost
down at large radii, validation-clean.
\end{tcolorbox}

\subsection{H8: The Equirect Is Fixed 2048$\times$1024 on Every Preset}
\label{sec:h8}

\textbf{Severity: High.}

\texttt{SkyRenderer::createOffscreenTargets} hard-codes
\texttt{offscreenWidth = 2048}, \texttt{offscreenHeight = 1024}
(\texttt{SkyRenderer.cpp:156--158}) regardless of output resolution or
graphics preset, then allocates three frame slots (24\,MiB) and shades
2.1\,M sky pixels per bake (\texttt{renderOffscreen}, \texttt{:240--330}).
On Minimal at 720p the bake shades more sky pixels than the screen has,
for a target sampled exactly once per background pixel through a bilinear
fetch (\texttt{postprocess.frag:114}) --- a fetch that cannot resolve detail
finer than the screen anyway. The C1 content-tiering precedent (report~23
C1, \texttt{Settings::textureArraySize}) applies directly: bake resolution
is a quality tier, not a constant.

\textit{Technique: tier the bake.} Scale the equirect ($1024\times512$ on
Minimal/low, $2048\times1024$ above) via the existing settings tier path;
recreate on tier change through the same resize flow that already rebuilds
these targets (\texttt{SceneRenderer.cpp:248--250}).

\begin{tcolorbox}[promptbox={Prompt --- H8: tiered equirect resolution}]
Goal: the bake costs what the screen can show. Steps: (1) add an equirect
resolution tier to \texttt{Settings} (default $2048\times1024$, low
$1024\times512$); (2) thread it through \texttt{createOffscreenTargets};
(3) recreate on tier change. Files: \texttt{utils/Settings.hpp},
\texttt{vulkan/renderer/SkyRenderer.cpp}. Validation: sky A/B per tier,
bake cost halved at low tier, validation-clean.
\end{tcolorbox}

% =====================================================================
\section{Medium-Severity Issues}
% =====================================================================

\subsection{M9: Sky Uploads Are Unconditional Per-Frame Host Writes}
\label{sec:m9}

\textbf{Severity: Medium.}

Two host writes fire every frame whether or not the sky changed:
\texttt{SkySphere::update} memcpys the whole \texttt{SkyUniform} from the
widget settings (\texttt{SkySphere.cpp:57--71}), called unconditionally from
the frame update (\texttt{MyApp.cpp:1167--1168}); and
\texttt{renderOffscreen} rewrites the shared UBO via
\texttt{updateUniformBuffer} on every bake (\texttt{SkyRenderer.cpp:250--252}).
Each is tens of bytes into coherent mapped memory --- noise-floor cost, but
also noise-floor fix: the equirect signature (\S\ref{sec:h4}) already
computes exactly the dirty bit these writes need. Reuse it (extended per
H4) so static-sky frames perform zero sky host writes.

\begin{tcolorbox}[promptbox={Prompt --- M9: dirty-gated sky uploads}]
Goal: a static sky costs zero host writes. Steps: (1) reuse the (fixed) sky
signature as the dirty bit for \texttt{SkySphere::update}; (2) skip the
\texttt{renderOffscreen} UBO rewrite when the signature matches; (3) force
dirty on settings-struct replacement. Files: \texttt{MyApp.cpp},
\texttt{vulkan/SkySphere.cpp}, \texttt{vulkan/renderer/SkyRenderer.cpp}.
Validation: byte-identical UBO contents, upload count zero on static frames,
validation-clean.
\end{tcolorbox}

\subsection{M10: The Composite Rebuilds Its UBO and Signature Every Frame}
\label{sec:m10}

\textbf{Severity: Medium.}

\texttt{PostProcessRenderer::render} maps, memcpys, and unmaps the
\texttt{WaterUBO} (\texttt{PostProcessRenderer.cpp:386--389}), rebuilds the
17-element \texttt{imageInfos} array (\texttt{:415--437}), and recomputes the
per-slot descriptor signature (\texttt{:455--467}) on every frame ---
including frames where \texttt{brushAlpha}/\texttt{brushMode} (set from the
selected brush entry, \texttt{MyApp.cpp:3179--3185}), matrices, and all
views are unchanged. The signature cache already makes the steady state
\textquotedblleft zero descriptor updates\textquotedblright\ (\texttt{:451--454});
gating the UBO memcpy and the signature rebuild on the same change dormancy
removes the remaining per-frame host work (map/unmap pair, 17-view array
fill, signature compare setup). Small alone; it sits on the hottest CPU
path in the frame.

\begin{tcolorbox}[promptbox={Prompt --- M10: change-gated composite upload}]
Goal: a static composite setup costs no host work. Steps: (1) cache the last
uploaded UBO bytes + view handles and skip map/memcpy/signature rebuild when
all match; (2) force refresh on resize/recreate. Files:
\texttt{vulkan/renderer/PostProcessRenderer.cpp}. Validation: byte-identical
UBO stream, zero uploads on static frames, validation-clean.
\end{tcolorbox}

\subsection{M11: Disabled Debug Legs Still Cost Fullscreen Fetches}
\label{sec:m11}

\textbf{Severity: Medium.}

The producers are gated --- the SDF task elides entirely when disabled
(\texttt{MyApp.cpp:2275--2300}), leaving cleared targets --- but the
consumer is not: the composite fetches \texttt{sdfDepthTex} and
\texttt{bboxDepthTex} unconditionally (\texttt{postprocess.frag:250,258}) and
tests \texttt{$< 1.0$} after the fetch, so every pixel of every frame pays
two depth fetches for overlays that are off. Same for vegetation depth when
vegetation is disabled (\texttt{:142}; the CPU already branches on
\texttt{settings.vegetationEnabled} for the cull at \texttt{MyApp.cpp}
cull task). The host knows all three flags; the shader does not. Pass them
through the UBO (three valid bits next to \texttt{waterBlurEnabled}, same
pattern as report~19 H4) and skip the fetches.

\begin{tcolorbox}[promptbox={Prompt --- M11: valid-bit the debug/veg fetches}]
Goal: a disabled overlay costs zero fetches. Steps: (1) add
\texttt{vegValid}/\texttt{sdfValid}/\texttt{bboxValid} to the composite UBO;
(2) skip each depth+color fetch when its bit is clear; (3) set the bits from
the existing settings at the \texttt{MyApp.cpp:3152--3225} call site. Files:
\texttt{shaders/postprocess.frag},
\texttt{vulkan/renderer/PostProcessRenderer.cpp},
\texttt{vulkan/ubo/WaterUBO.hpp}, \texttt{MyApp.cpp}. Validation:
pixel-identical output in all flag combinations, fetch count down when off,
validation-clean.
\end{tcolorbox}

\subsection{M12: Tail Targets Are Full-Res, Full-Slot, and Resize-Churned}
\label{sec:m12}

\textbf{Severity: Medium.}

\S\ref{sec:vram}: $\sim$95\,MiB at 1080p, $\sim$380\,MiB at 4K, all
resolution-scaled and none content-scaled. Three compounding details: brush
color+depth+backface allocate at full render resolution though brush editing
happens at stroke scale (a half-res brush tail would composite indistinguishably
under 0.5 alpha); all targets allocate all 3 slots even when the producing
pass is elided for the session (C1, SDF-style); and every swapchain resize
destroys and recreates the brush targets from scratch
(\texttt{BrushRenderer.cpp:104--115}) plus the backface set, with layout
tracking reset to \texttt{UNDEFINED} --- no preservation, no lazy
re-allocation. At 4K this section alone commits more than report~23's entire
vegetation+billboard+impostor tally.

\textit{Technique: scale, share, preserve.} Render the brush tail at a
brush-quality scale factor (settings-tiered, default 1.0, low 0.5) with the
composite upsampling the brush color (depth tests stay conservative via the
closest-tap rule already used for water/veg); allocate backface depth only
while PAINT mode is armed; preserve images across resizes when the size is
unchanged.

\begin{tcolorbox}[promptbox={Prompt --- M12: scaled, lazy tail targets}]
Goal: the tail commits what editing needs. Steps: (1) add a brush-tail scale
to \texttt{Settings} and size brush targets by it; (2) allocate backface
depth lazily on PAINT-mode arming; (3) skip recreate when extents match.
Files: \texttt{utils/Settings.hpp},
\texttt{vulkan/renderer/BrushRenderer.cpp,.hpp},
\texttt{vulkan/renderer/BrushBackFaceRenderer.cpp}. Validation: overlay A/B
per tier, VRAM delta per tier, validation-clean.
\end{tcolorbox}

% =====================================================================
\section{Low-Severity / Hygiene Issues}
% =====================================================================

\subsection{L13: Two Fullscreen Vertex Shaders for Three Vertices}

\texttt{sky\_fullscreen.vert.spv} (2\,515 words / 515 instrs) vs
\texttt{fullscreen.vert.spv} (306 words / 64 instrs): the sky variant
reconstructs world position through the invVP matvec and exports two
varyings (\texttt{shaders/sky\_fullscreen.vert}) to shade 3 vertices, while
the composite variant emits bare clip positions
(\texttt{shaders/fullscreen.vert}). The reconstruct exists to feed
\texttt{sky.frag}'s direction-only use (\texttt{fragPosWorld} minus camera
position, normalized) --- a frag-side reconstruct from \texttt{gl\_FragCoord}
and a push-constant invVP would let one shared triangle vertex shader serve
sky, composite, and equirect. Trivial ALU either way (3 vertices); the real
cost is pipeline/shader-module sprawl: three fullscreen pipelines with
near-identical vertex stages.

\subsection{L14: Composite Viewport Follows Cached State, Not the Swapchain}

\texttt{PostProcessRenderer::render} sets viewport/scissor from
\texttt{renderWidth/Height} (\texttt{PostProcessRenderer.cpp:577--590}),
cached via \texttt{setRenderSize} (\texttt{:45--48}) and fed from app
dimensions at init and resize (\texttt{SceneRenderer.cpp:246,885}). Any
drift between those dimensions and the actual swapchain extent silently
scales or clips the one pass that defines the final image --- and nothing
asserts their equality. Derive the viewport from the swapchain extent at
record time (it is already on hand where the composite is recorded,
\texttt{MyApp.cpp:3152--3225}) or debug-assert equality after resize.

\subsection{L15: One-Shot DIAG Block and Comment Drift in the Tail}

\texttt{PostProcessRenderer.cpp:393--414} prints a \texttt{static
bool}-gated DIAG line of raw view handles to stderr on the first composite
frame of every run --- the report-22-era debug leftover pattern L-sections
keep finding. Nearby drift: the \texttt{brushMode} float documents only
\texttt{0=overlay, 2=PAINT} (\texttt{postprocess.frag:23}) while the brush
manager defines the full mode set; the 17-binding layout header comment
(\texttt{PostProcessRenderer.cpp:61}) survives each binding repurposing
(M11's valid bits will be next). Delete the DIAG, document the mode enum in
one place, and keep the binding table truthful.

% =====================================================================
\section{Implementation Roadmap and Expected Deltas}
% =====================================================================

\begin{table}[h]
\centering
\small
\begin{tabular}{@{}lp{5.8cm}p{5.6cm}@{}}
\toprule
\textbf{Item} & \textbf{Change} & \textbf{Est.\ effect} \\
\midrule
C1 & Content-gated brush tail & $-3$ full-res passes + 1 submit/frame when idle \\
C2 & Conditional sky FS + day-factored stars & $-2$M sky px/frame when solid off; $-$2 hashes/px by day \\
C3 & Lazy composite prologue + valid flags & $-$6 fetches + invVP + atan/acos per solid px \\
H4 & Complete equirect signature & stale-sky class of bug closed \\
H5 & De-tessellated brush depth & HS/DS off the brush depth path \\
H6 & Cheap brush FS & brush ALU down $\sim$10$\times$ \\
H7 & Bounded water-blur kernel & $-$sparse taps at large radii \\
H8 & Tiered equirect & $-1$M bake px + $-18$\,MiB at low tier \\
M9--M12 & Dirty-gated uploads, fetch valid bits, scaled targets & additive host + VRAM \\
L13--L15 & Shared triangle VS, viewport assert, DIAG/comments & noise floor \\
\bottomrule
\end{tabular}
\end{table}

C1--C3 are the items that change what the engine \emph{costs per frame} on
the tail: idle submit/pass overhead, unconditional sky shading, and the
per-pixel prologue tax. H4 is a correctness fix with a performance
signature. None of the changes touches scene content except H6's brush
shading (an editing overlay) and H7--H8's blur/bake tiers (quality knobs).

Explicitly out of scope (audited shallowly or not at all; candidates for
report~25): the impostor bake's synchronous 20-view capture
(\texttt{ImpostorCapture.cpp:240--336}, one blocking submit with 20
\texttt{BeginRendering}+indexed draws, 40 \texttt{ImGui\_RemoveTexture} per
capture); the GDAL heightmap load path (\texttt{math/HeightMapTif.cpp},
blocking \texttt{GDALOpen} + whole-image double-buffer +
non-contiguous \texttt{vector<vector<float>>}); CPU threading
(\texttt{space/ThreadPool.cpp}, single-mutex FIFO, no work-stealing);
input/events (\texttt{events/EventManager.cpp}, per-publish snapshot +
double-loop drain); and the ImGui widget layer (release-only, never
profiled).

% =====================================================================
\section{Overall Score}
\label{sec:score24}
% =====================================================================

\subsection{Rubric (Report-24 Scope)}

\begin{table}[h]
\centering
\small
\begin{tabular}{@{}lcp{7.2cm}@{}}
\toprule
\textbf{Category} & \textbf{Wt.} & \textbf{Score and justification} \\
\midrule
Composite efficiency & 35\% & \textbf{5.0} --- blur gating, LOD discipline,
descriptor signature cache and timeline waits are good; the unconditional
prologue (invVP + atan/acos + sky + 4 extra fetches per solid pixel) drags
it down. \\
Sky & 30\% & \textbf{5.5} --- the equirect cache removes 2\,M px/frame when
static and the fullscreen triangle replaced sphere geometry; the draw is
unconditional, day pixels hash stars for nothing, and the signature misses
five inputs. \\
Brush tail & 25\% & \textbf{4.5} --- dedicated IRs, own queue, and
layout-tracked targets are good; three full-res passes + submit ungated,
tessellated depth-only draws, and full-recipe shading are not. \\
Upload / target hygiene & 10\% & \textbf{6.0} --- coherent mapped uploads,
no per-frame allocations in the tail; uploads and 95\,MiB of targets ignore
change and content. \\
\bottomrule
\end{tabular}
\end{table}

\subsection{Overall: 5.1/10 Scope, 6.1/10 Combined}

\begin{equation}
0.35(5.0) + 0.30(5.5) + 0.25(4.5) + 0.10(6.0) = \mathbf{5.1}
\end{equation}

Report~23 scored the combined engine 6.2/10; that estimate priced geometry,
submits, pools, arrays, mixer, and load, but not the per-pixel tail.
Measuring it:

\begin{equation}
0.90(6.2) + 0.10(5.1) = \mathbf{6.1}
\end{equation}

The engine's verdict is unchanged in kind, extended in scope once more: the
frame graph architecture (async queues, timeline DAG, cached descriptors,
cached equirect, elided debug tasks) is sound; the constants --- which
passes run unconditionally, and what each pixel computes before its first
branch --- are again where the money is. Landing C1--C3 moves this scope
toward $\sim$8/10 with no redesign.

\end{document}
